\begin{theorem}[Markov’s Inequality]
Let \((\Omega,\mathcal F,\mathbb P)\) be a probability space and let
\(X:\Omega\to[0,\infty]\) be a non‑negative random variable with finite
expectation \(\mathbb E[X]\). For any constant \(a>0\),
\[
\mathbb P\!\bigl(X\ge a\bigr)\le\frac{\mathbb E[X]}{a}.
\]
Equality holds iff \(X=a\,\mathbf 1_{\{X\ge a\}}\) a.s. (i.e. \(X\) takes only
the two values \(0\) and \(a\)).
\end{theorem}

\begin{proof}
\textbf{Assumptions and edge‑case analysis.}  We assume \(X(\omega)\ge0\) for
every \(\omega\in\Omega\) and \(a>0\).  If \(\mathbb P(X\ge a)=0\) the
inequality is immediate.  If \(\mathbb E[X]=\infty\) the right‑hand side is
\(+\infty\) and the inequality holds trivially.  The case \(a=0\) is
excluded by hypothesis.

\medskip
\textbf{Indicator of the event \(\{X\ge a\}\).}  Define the random variable
\[
\mathbf 1_{\{X\ge a\}}(\omega)=
\begin{cases}
1, & \text{if }X(\omega)\ge a,\\[2mm]
0, & \text{otherwise}.
\end{cases}
\]
By the definition of expectation of an indicator,
\begin{equation*}
\mathbb E\!\bigl[\mathbf 1_{\{X\ge a\}}\bigr]=\mathbb P\!\bigl(X\ge a\bigr).
\tag{2.1}
\end{equation*}

\medskip
\textbf{Pointwise domination.}  For each \(\omega\in\Omega\) we have
\begin{equation*}
X(\omega)\;\ge\;a\,\mathbf 1_{\{X\ge a\}}(\omega). \tag{3.1}
\end{equation*}
\emph{Verification.}  If \(X(\omega)\ge a\) then
\(\mathbf 1_{\{X\ge a\}}(\omega)=1\) and the right‑hand side equals \(a\le
X(\omega)\).  If \(X(\omega)<a\) then the indicator is zero and the
inequality reduces to \(X(\omega)\ge0\), which holds by non‑negativity of
\(X\).

\medskip
\textbf{Taking expectations and applying monotonicity.}  Since (3.1) holds
pointwise, monotonicity of expectation yields
\begin{equation*}
\mathbb E[X]\;\ge\;\mathbb E\!\bigl[a\,\mathbf 1_{\{X\ge a\}}\bigr].
\tag{4.1}
\end{equation*}
Linearity of expectation gives
\begin{equation*}
\mathbb E\!\bigl[a\,\mathbf 1_{\{X\ge a\}}\bigr]
   =a\,\mathbb E\!\bigl[\mathbf 1_{\{X\ge a\}}\bigr].
\tag{4.2}
\end{equation*}
Substituting (2.1) into (4.2) and then into (4.1) we obtain
\begin{equation*}
\mathbb E[X]\;\ge\;a\,\mathbb P\!\bigl(X\ge a\bigr).
\tag{4.3}
\end{equation*}

\medskip
\textbf{Division by the positive constant \(a\).}  Because \(a>0\),
\begin{equation*}
\mathbb P\!\bigl(X\ge a\bigr)\;\le\;\frac{\mathbb E[X]}{a}.
\tag{5.1}
\end{equation*}
This is precisely Markov’s inequality.

\medskip
\textbf{Equality condition.}  Equality in (5.1) can occur only if every
inequality used above is an equality.

Equality in the monotonicity step (4.1) requires that the pointwise
inequality (3.1) be an equality \(\mathbb P\)-almost everywhere, i.e.
\begin{equation*}
X(\omega)=a\,\mathbf 1_{\{X\ge a\}}(\omega)\qquad\text{a.s.}
\tag{6.1}
\end{equation*}
Thus \(X\) takes only the values \(0\) and \(a\) almost surely.  Conversely,
if \(X\) assumes only the two values \(0\) and \(a\) with probabilities
\(p_{0}\) and \(p_{a}\) (\(p_{0}+p_{a}=1\)), then
\(\mathbb P(X\ge a)=p_{a}\) and \(\mathbb E[X]=a\,p_{a}\), so (5.1) holds
with equality.

Hence equality holds \emph{iff}
\[
X=a\,\mathbf 1_{\{X\ge a\}}\quad\text{a.s.},
\]
i.e. the random variable takes only the values \(0\) and \(a\). \(\square\)
\end{proof}